\section{Orbit-Breaking Candidates}

The experiments considered two main orbit-breaking fixed-length candidates.

\paragraph{Block heat-bath.}
Choose a contiguous block $B=(t_i,\ldots,t_{i+r-1})$, compute its product
$h=t_i\cdots t_{i+r-1}$, and replace $B$ by an exactly uniform element of
\[
  \mathcal{B}_r(h)
  =
  \{(s_1,\ldots,s_r)\in T_N^r : s_1\cdots s_r=h\}.
\]
The outside coordinates are unchanged.

\paragraph{Square refresh plus Hurwitz scrambling.}
If an adjacent block has the form $(t,t)$, replace it by $(s,s)$ for a random
transposition $s$, then apply Hurwitz scrambling.  This move preserves product
and length.  It can break Hurwitz components, but later measurements showed
that a random-slot square refresh acts too rarely at moderate $N$.

\subsection{Stationarity of Exact Block Heat-Bath}

\begin{proposition}[Reversibility of block heat-bath]
Fix a block-size distribution and suppose each selected block is resampled
exactly uniformly from $\mathcal{B}_r(h)$, where $h$ is the product of the old
block.  The resulting random-scan block heat-bath chain is reversible with
respect to the uniform distribution on $\Omega(N,m)$.
\end{proposition}

\begin{proof}
Condition on the choice of block position and size, and fix the outside
coordinates.  The outside coordinates force the block product $h$ because the
total word product must remain $\id$.  Under the uniform law on $\Omega(N,m)$,
the conditional distribution of the selected block is precisely the uniform
law on $\mathcal{B}_r(h)$.  Replacing the block by an independent sample from
that conditional distribution is the usual Gibbs update, hence it satisfies
detailed balance.  Averaging over block choices preserves reversibility.
\end{proof}

Consequently, if a lazy version of the chain is irreducible on $\Omega(N,m)$,
it converges to the global uniform fixed-length bridge distribution.  The
remaining tasks are therefore:
\begin{enumerate}
\item sample $\mathcal{B}_r(h)$ exactly and efficiently for useful $r$;
\item prove irreducibility for the selected block sizes;
\item prove or estimate a useful mixing time.
\end{enumerate}
